\section{Towards better Multi-Agent LLM Systems}
\label{sec:discussions}

\begin{figure*}[ht]
\centering
\includegraphics[width=0.95\linewidth]{figures/masft_bar.pdf}
\caption{Distribution of failure in \dataset{} with \taxonomy{} labels on total 210 traces. This plot visualizes the failure distributions of the first 30 traces for each system. As the specific tasks and benchmarks may differ across the MAS configurations shown, these results are intended to illustrate system-specific failure profiles rather than to serve as a performance comparison across MAS.}
\label{fig:mas_failure_bar}
\end{figure*}

We now discuss the broader implications and usage of \dataset{} and \taxonomy{}. \dataset{}, with its annotations grounded in \taxonomy{}, provides crucial empirical evidence, while \taxonomy{} offers a foundational framework and practical tool for understanding, debugging, and ultimately improving MAS. By concretely defining failure modes and providing a large-scale dataset of their occurrences, our work outlines the challenges in building reliable MAS and opens targeted research problems for the community.
% This section highlights how \dataset{} and \taxonomy{} together aid agentic system development, suggesting that progress requires focusing on system design alongside model capabilities.

\subsection{Failure Breakdown in \dataset{}}
Our analysis of \dataset{} reveals that failure distributions differ markedly across various MAS, often reflecting their unique architectural characteristics and design philosophies. For example, as illustrated in Figure~\ref{fig:mas_failure_bar}, we observe specific patterns: AppWorld frequently suffers from premature terminations (FM-3.1), potentially due to its star topology and lack of a predefined workflow making termination conditions less obvious; OpenManus exhibits a tendency towards step repetition (FM-1.3); and HyperAgent could benefit from addressing its dominant failure modes of step repetition (FM-1.3) and incorrect verification (FM-3.3). These system-specific profiles underscore that there is no one-size-fits-all solution to MAS failures.

We also use \dataset{} to study the impact of different underlying language models and MAS designs on failure patterns. For instance, when comparing GPT-4o and Claude 3.7 Sonnet within the MetaGPT framework on programming tasks, we find that while GPT-4o generally performs better than Claude, it shows significantly fewer FC1 (System Design Issues) failures by 39\%. We also examine the impact of different MAS designs on the same benchmark, such as comparing MetaGPT and ChatDev on ProgramDev. Here, while MetaGPT generally outperforms ChatDev by having 60-68\% less failure in FC1 and FC2, it has 1.56x more FC3 failure than ChatDev. These comparative analyses, detailed further in Appendix~\ref{sec:failure_across_model}, provide insights into how model choice and architectural patterns influence system performance and distribution of failures.

\subsection{\taxonomy{} as a Practical Development Tool}
\label{sec:mast-utility}

Developing robust MAS is challenging: aggregate success rates can obscure the specific impacts of optimizations. \taxonomy{} addresses this by providing a structured vocabulary for systematic failure breakdown. Using our LLM annotator with \taxonomy{}, developers can obtain quantitative analyses of failure profiles for specific systems. We demonstrate \taxonomy{}'s practical usage in guiding MAS improvement in our case studies (Appendix~\ref{sec:case-studies}). The Failure Mode breakdown analysis (Appendix~\ref{sec:interventions_effects_on_failure_modes_with_mast}) shows which failure modes were mitigated and reveals any resulting trade-offs. This granular view, moving beyond aggregate metrics, is crucial for understanding \textit{why} an intervention works and for iterating effectively towards more robust systems.

% Developing robust MAS presents significant challenges. Pinpointing causes for MAS failures is difficult without a systematic framework, especially with varied failure manifestations. Consequently, developers often resort to ad-hoc debugging \cite{fritzson1992generalized}, and evaluating optimization impacts is complex, as aggregate success rates can obscure specific effects. \taxonomy{} addresses this by offering a structured framework for systematic diagnosis. Using our LLM annotator with \taxonomy{}, developers can obtain quantitative failure breakdowns from MAS execution traces to guide debugging.

% To illustrate, we conduct two small-scale case studies (Appendix~\ref{sec:case-studies}) on the effects of prompt and architectural solutions on MAS, where \taxonomy{} facilitates a detailed understanding of intervention impacts. Instead of relying solely on aggregate success rates, developers can perform before-and-after comparisons using \taxonomy{}. Our case studies on ChatDev and AG2 (Table~\ref{tab:ag2-chatdev-accuracies}) demonstrate that a \taxonomy{}-based analysis (detailed in Appendix~\ref{sec:interventions_effects_on_failure_modes_with_mast}) reveals which specific failure modes were mitigated and whether any trade-offs occurred. This detailed view is crucial for understanding \textit{why} an intervention works and for iterating effectively towards more robust systems.

% Developing robust MAS presents significant challenges. When a system fails frequently on a benchmark (e.g., 66.7\% failure for ChatDev on ProgramDev, Figure~\ref{fig:mas_failure_rate}), pinpointing underlying causes is difficult, particularly if failure manifestations vary widely. Without a systematic framework like \taxonomy{}, developers often resort to ad-hoc debugging of individual failed traces \cite{fritzson1992generalized}. Furthermore, evaluating interventions is complex; modest gains in overall success rates might obscure whether a fix addressed intended issues or introduced new problems.

% \taxonomy{} offers practical value by enabling the detection and breakdown of failures. Its structured vocabulary and clear definitions for distinct failure modes facilitate systematic diagnosis. When our LLM annotator is used to analyze MAS execution traces based on \taxonomy{}, developers can obtain a quantitative breakdown of failure types occurring in their system, pinpointing frequent failure modes and guiding debugging efforts toward high-impact areas.

% To illustrate, we conduct two small-scale case studies (Appendix~\ref{sec:case-studies}) on the effects of prompt and architectural solutions on MAS. \taxonomy{} facilitates a detailed understanding of the impact of these interventions. Instead of relying solely on aggregate success rates, we can perform before-and-after comparisons using \taxonomy{}. Our case studies, involving interventions on ChatDev and AG2 (Table~\ref{tab:ag2-chatdev-accuracies}), show that a \taxonomy{}-based analysis (detailed in Appendix~\ref{sec:interventions_effects_on_failure_modes_with_mast}) reveals which specific failure modes were mitigated and whether any trade-offs occurred. This detailed view guides the understanding of \textit{ why} an intervention works and the effective iteration toward more robust systems.

\subsection{Beyond Model Capabilities: The Primacy of System Design}
\label{sec:system-design-matters}
While one could simply attribute failures in \dataset{} to limitations of present-day LLM (e.g., hallucinations, misalignment), we conjecture that improvements in the base model capabilities will be insufficient to address the full \taxnomony.  
Instead, we argue that good MAS design requires organizational understanding -- even organizations of sophisticated individuals can fail catastrophically~\citep{perrow1984normal} if the organization structure is flawed. Previous research in high-reliability organizations has shown that well-defined design principles can prevent such failures \cite{roberts1989high, ctx26792789720006532}. 

Consistent with organization theories, our findings indicate that many MAS failures arise from the challenges in organizational design and agent coordination rather than the limitations of individual agents. In our intervention case studies (Appendix~\ref{sec:case-studies}), we apply MAS system workflow and prompt changes respectively (results in Table~\ref{tab:ag2-chatdev-accuracies}). With the same underlying model, we achieve max improvements of 15.6\%. This highlights that MAS failures can be address with better system designs.

Although first step interventions lead to performance gains, not all failure modes are resolved, and task completion rates still remain low, indicating that more substantial improvements are needed. Achieving high reliability may requires combinatorial changes ranging from agent system organization to model level improvements (see Table~\ref{tab:solution_vs_failure_modes}). \taxonomy{}, by providing a clear framework of failure points identified from \dataset{}, helps identify where these structural weaknesses lie and can guide the design and evaluation of more sophisticated MAS architectures.

%%%%%%%%%%